% !TeX program = xelatex
% Regression for Persian punctuation, LTR technical terms and a C listing.
\documentclass[12pt,a4paper]{article}
\usepackage{tex/persian-os}
\renewcommand{\OSFaculty}{دانشکده نمونه}
\renewcommand{\OSCourse}{درس آزمایشی}
\renewcommand{\OSAssignment}{تمرین نمونه}
\renewcommand{\OSTerm}{نیمسال اول \OSDateRange{۱۴۰۵}{۱۴۰۶}}
\renewcommand{\OSQuestionNumber}{۲}
\begin{document}
\OSMakeQuestionTitle
\OSAnswerHeading
\OSPart{آ}[نشانه‌گذاری و اصطلاحات]
«پردازه» آماده است، اما فرایند دیگری در حال اجراست؛ آیا \OSLatin{fork()} فراخوانی شده است؟

\OSTechnicalTerm{فراخوانی سیستمی}{waitpid()} در این مسئله استفاده می‌شود.

این عبارت \OSParen{آزمون فارسی} است؛ همچنین \OSBracketed{نوشتهٔ فارسی} به کار می‌رود.

در یک جملهٔ بلند، \OSBracketed{یک عبارت فارسی طولانی در چند سطر که با نقطه، ویرگول، علامت سؤال؟ و کلمهٔ انگلیسی \OSLatin{fork()} ترکیب می‌شود و باید همچنان کروشه‌های ابتدا و انتهای آن به سمت داخل جمله قرار بگیرند} را نیز می‌آزماییم.
\OSPart{ب}[نمونهٔ کد]
نمونهٔ برنامهٔ زبان سی به‌صورت چپ‌به‌راست نمایش داده می‌شود:
\begin{latin}
\begin{lstlisting}[style=OSCode]
#include <stdio.h>
int main(void) {
    printf("Hello, OS!\\n");
    return 0;
}
\end{lstlisting}
\end{latin}
\end{document}
